\section{Hierarchical Coordination of Large-Scale Fleets under Temporal Tasks}
\label{sec:hierarchical_coordination}

Large-scale temporal-task coordination for heterogeneous fleets faces two
coupled sources of complexity. The first source is logical. Temporal missions
must be decomposed into admissible task sequences that satisfy ordering,
concurrency, and completion requirements. The second source is combinatorial.
Concrete robots with different capabilities must be grouped into teams and
dispatched over the workspace with bounded response time. Direct formulations
that synchronize temporal automata with all robot models can provide strong
correctness guarantees, but their size grows rapidly with the fleet cardinality
and the mission complexity
\cite{schillinger2018simultaneous,kantaros2020stylus,luo2021temporal,
leahy2021scalable}. Recent advances have improved scalability by sampling,
decision-tree expansion, partial-order reduction, and hierarchical task
abstractions
\cite{chen2025real,liu2024time,luo2025simultaneous,gosrich2025online,
luo2025hulk}. Nevertheless, large fleets still require a sharper separation
between temporal decomposition, team formation, and execution-level
realization.

This section develops such a separation. The emphasis is the coordination
backbone that connects temporal specifications to executable team plans. The
presentation deliberately suppresses interface-level supervision and keeps
re-optimization details to a minimum. Those aspects are deferred to the next
section. The resulting architecture contains three tightly coupled layers. The
fleet level reasons over temporal tasks and constructs abstract team plans. The
intermediate level instantiates these abstract plans by assigning concrete
robots to capability profiles. The execution level converts each team plan into
robot-level plans $\tau_i$, which have been defined previously as timed action
sequences.

\subsection{Hierarchical coordination model}
\label{subsec:hierarchical_model}

Let $\mathcal{N}\triangleq\{1,\cdots,N\}$ denote the fleet and let
$\Phi_t\triangleq\{\varphi_1,\cdots,\varphi_M\}$ denote the set of temporal
missions known at time $t$. Each mission $\varphi_\ell$ is interpreted as a
syntactically co-safe temporal formula over collaborative tasks. This modeling
choice is natural for finite-horizon missions in which completion can be
certified after a finite trace, and it is consistent with automata-based
planning for temporal logic \cite{baier2008principles}. The task abstraction
must capture both capability requirements and the spatial support over which
collaboration is executed.

\begin{definition}[Collaborative task]
\label{def:collaborative_task}
A collaborative task is a tuple
\[
\omega \triangleq \big(S_\omega,\eta_\omega,\mathcal{J}_\omega\big),
\]
where $S_\omega\subseteq\mathcal{W}$ is the execution region,
$\eta_\omega$ is a duration map, and
$\mathcal{J}_\omega\triangleq\{(n_j,a_j,s_j)\}_{j=1}^{J_\omega}$ is the set
of subtasks. For each subtask, $n_j$ is the required number of robots,
$a_j$ is the required action, and $s_j\in S_\omega$ is the service location.
\end{definition}

Definition~\ref{def:collaborative_task} extends purely counting-based task
descriptions by retaining the spatial support and the dependence of execution
time on the assigned team. This distinction is important in large-scale fleets,
because the same cardinality requirement may induce very different response
times when the available robots differ in location and capability
\cite{sahin2019multirobot,cardona2024planning}. It is also convenient for
hierarchical coordination, since the fleet level only needs abstract task
information, while the execution level resolves robot trajectories.

The intermediate objects of the hierarchy are task sequences assigned to teams
and the associated capacity profiles.

\begin{definition}[Abstract team plan and capacity profile]
\label{def:team_profile}
For team index $k\in\mathcal{K}$, an abstract team plan is a finite sequence
\[
\Gamma_k \triangleq \omega_k^1\omega_k^2\cdots\omega_k^{L_k}.
\]
The capacity profile induced by $\Gamma_k$ is the map
$\beta_k:\mathcal{A}\rightarrow\mathbb{N}$ defined by
\begin{equation}
\label{eq:cap_profile_ch4}
\beta_k(a) \triangleq
\max_{\omega\in\Gamma_k}
\max_{\substack{(n_j,a_j,s_j)\in\mathcal{J}_\omega\\ a_j=a}}
n_j,
\end{equation}
where $\beta_k(a)$ is taken as zero when action $a$ does not appear in
$\Gamma_k$.
\end{definition}

Equation~\eqref{eq:cap_profile_ch4} records the largest instantaneous demand
for each action along the team sequence. It deliberately does not identify
which robots realize that demand. This abstraction is the key mechanism that
decouples fleet-level temporal decomposition from concrete robot assignment.

\subsection{Automaton-guided fleet-level assignment}
\label{subsec:automaton_assignment}

For each mission $\varphi_\ell\in\Phi_t$, let
\[
\mathcal{B}_{\varphi_\ell}
\triangleq
\big(Q_\ell,Q_\ell^0,\Sigma_\ell,\delta_\ell,F_\ell\big)
\]
denote the corresponding nondeterministic B\"uchi or finite-word acceptor,
depending on the exact temporal fragment introduced previously. The fleet-level
coordination problem is to construct a set of abstract team plans
$\{\Gamma_k\}_{k\in\mathcal{K}}$ and an admissible number of teams $K$ such
that the induced trace satisfies all mission automata and the aggregate
capability demand remains feasible.

\begin{problem}
\label{prob:fleet_level_coordination}
Given the mission set $\Phi_t$ and the fleet capability model
$\{\mathcal{A}_i\}_{i\in\mathcal{N}}$, determine the number of teams $K$, the
abstract team plans $\{\Gamma_k\}_{k\in\mathcal{K}}$, and the associated
capacity profiles $\{\beta_k\}_{k\in\mathcal{K}}$ such that all missions are
accepted by their automata and the fleet response time is minimized.
\end{problem}

A direct synchronization of all robot models with all mission automata is
avoided. Instead, the search is carried out over partial team plans together
with the currently reachable automaton states.

\begin{definition}[Assignment node]
\label{def:assignment_node}
Let $\mathcal{M}\triangleq\{1,\cdots,M\}$. An assignment node is a tuple
\[
\nu \triangleq
\big(\{\Gamma_k\}_{k\in\mathcal{K}},\{\widehat{Q}_\ell\}_{\ell\in\mathcal{M}}\big),
\]
where $\Gamma_k$ is the current task sequence of team $k$ and
$\widehat{Q}_\ell\subseteq Q_\ell$ is the set of automaton states reachable in
$\mathcal{B}_{\varphi_\ell}$ after executing the task symbols already assigned
in the node.
\end{definition}

The root node contains empty team plans and the initial reachable-state sets.
Each expansion appends one enabled task symbol to one existing team, or to a
newly created team, and updates the reachable-state sets accordingly. A task
symbol is enabled at node $\nu$ if it appears on at least one outgoing
transition of some state in $\widehat{Q}_\ell$ for at least one mission
$\ell\in\mathcal{M}$. This rule allows temporal decomposition and team
assignment to proceed simultaneously.

To rank candidate nodes, a multi-term value function is introduced:
\begin{equation}
\label{eq:node_value_ch4}
\chi(\nu) \triangleq
\max_{k\in\mathcal{K}} T_k
+ \lambda_c \sum_{k\in\mathcal{K}} C_k
+ \lambda_p \sum_{\ell\in\mathcal{M}}
w_\ell \min_{q\in\widehat{Q}_\ell}\psi_\ell(q,F_\ell),
\end{equation}
where $T_k$ is the estimated completion time of $\Gamma_k$, $C_k$ is the
estimated travel-and-service cost of $\Gamma_k$, $w_\ell$ is the prescribed
importance weight of mission $\varphi_\ell$, and
$\psi_\ell(q,F_\ell)$ is the shortest graph distance from $q$ to the accepting
set $F_\ell$,
where $\lambda_c>0$ and $\lambda_p>0$ are regularization weights.

Equation~\eqref{eq:node_value_ch4} has a clear interpretation. The first term
controls the makespan of the current partial assignment. The second term
penalizes unnecessarily expensive team sequences. The third term favors nodes
that are logically closer to acceptance. A purely makespan-driven expansion is
therefore avoided, which improves search stability when many nodes share the
same temporal depth.

Every node must also satisfy a fleet-level capability test. Using the capacity
profiles defined in Definition~\ref{def:team_profile}, feasibility requires
\begin{equation}
\label{eq:global_capacity_ch4}
\sum_{k\in\mathcal{K}} \beta_k(a)
\leq
\sum_{i\in\mathcal{N}}\mathds{1}\big(a\in\mathcal{A}_i\big),
\quad \forall a\in\mathcal{A},
\end{equation}
where the left-hand side is the total demand for action $a$ induced by the
current abstract team plans, and the right-hand side is the number of robots in
the fleet that can perform action $a$.

A further reduction is obtained through dominance pruning. For each node,
define the profile
\begin{equation}
\label{eq:node_profile_ch4}
\rho(\nu) \triangleq
\Big[
\{T_k\}_{k\in\mathcal{K}},
\{C_k\}_{k\in\mathcal{K}},
\Big\{
\min_{q\in\widehat{Q}_\ell}\psi_\ell(q,F_\ell)
\Big\}_{\ell\in\mathcal{M}}
\Big].
\end{equation}
Node $\nu_1$ dominates node $\nu_2$ when
$\rho(\nu_1)\leq\rho(\nu_2)$ componentwise and at least one inequality is
strict. Only non-dominated nodes are retained on the frontier. The search thus
progresses over a compressed representation of the assignment space instead of a
full synchronous product.

\begin{algorithm}[t]
\caption{Automaton-guided Hierarchical Task Assignment}
\label{alg:hierarchical_assignment}
\SetAlgoLined
\KwIn{Mission set $\Phi_t$, automata $\{\mathcal{B}_{\varphi_\ell}\}$,
fleet $\mathcal{N}$, horizon $H$.}
\KwOut{Number of teams $K$, abstract team plans $\{\Gamma_k\}$,
capacity profiles $\{\beta_k\}$.}

Initialize the root node
$\nu_0 \gets (\{\emptyset\},\{Q_\ell^0\}_{\ell\in\mathcal{M}})$\;
Initialize the frontier $\mathscr{F}\gets\{\nu_0\}$\;

\While{the budget is not exhausted}{
  Select a batch of frontier nodes with smallest
  $\chi(\nu)$ in~\eqref{eq:node_value_ch4}\;

  \ForEach{selected node $\nu$}{
    compute the enabled task-symbol set from
    $\{\widehat{Q}_\ell\}_{\ell\in\mathcal{M}}$\;

    \ForEach{enabled symbol $\omega$}{
      \ForEach{existing team $k$ and the option of creating one new team}{
        append $\omega$ to $\Gamma_k$\;
        update the reachable-state sets
        $\{\widehat{Q}_\ell\}_{\ell\in\mathcal{M}}$\;
        update $T_k$, $C_k$, and $\beta_k(\cdot)$\;
        \lIf{\eqref{eq:global_capacity_ch4} is violated}{discard the node}
        \lIf{the number of newly committed tasks exceeds $H$}{pause expansion}
      }
    }
  }

  prune dominated nodes using~\eqref{eq:node_profile_ch4}\;
  update the frontier with the non-dominated surviving nodes\;
}

return the complete frontier node with minimum
$\chi(\nu)$ in~\eqref{eq:node_value_ch4}\;
\end{algorithm}

Algorithm~\ref{alg:hierarchical_assignment} summarizes the fleet-level
backbone. The role of the horizon $H$ is purely computational in the present
section. It bounds the number of future task commitments in one optimization
cycle and prevents the assignment layer from becoming excessively myopic or
excessively rigid. Detailed event-triggered replanning policies are omitted
here and are addressed separately later.

The logical correctness of the fleet-level layer follows directly from the node
semantics and the acceptance condition of the mission automata.

\begin{proposition}
\label{prop:fleet_correctness}
Assume that the search horizon contains all admissible task symbols for the
current planning instance. If Algorithm~\ref{alg:hierarchical_assignment}
returns a complete feasible node, then the resulting abstract team plans
$\{\Gamma_k\}_{k\in\mathcal{K}}$ satisfy all mission formulas and all aggregate
capability constraints in~\eqref{eq:global_capacity_ch4}.
\end{proposition}

The proof is immediate. Each node expansion follows valid automaton
transitions, and completeness requires that every mission automaton reaches an
accepting set. Capability feasibility is enforced explicitly through
\eqref{eq:global_capacity_ch4}. Hence any complete feasible node corresponds to
a temporally valid fleet-level decomposition.

\subsection{Capacity-aware team formation}
\label{subsec:capacity_team_formation}

The fleet-level layer determines what each team must be able to do, but it does
not determine which robots belong to which team. The next layer instantiates
each abstract team plan by assigning concrete robots to the induced capacity
profile. This assignment must preserve disjoint team membership and may include
limited redundancy to absorb uncertainty in workload or travel time.

\begin{problem}
\label{prob:team_formation}
Given the abstract team plans $\{\Gamma_k\}_{k\in\mathcal{K}}$ returned by
Algorithm~\ref{alg:hierarchical_assignment}, determine the concrete teams
$\{\mathcal{N}_k\}_{k\in\mathcal{K}}$ such that each team can realize the
capacity profile induced by its plan and the worst team completion time is
minimized.
\end{problem}

Introduce binary variables $b_{ik}\in\{0,1\}$, where $b_{ik}=1$ means that
robot $i$ is assigned to team $k$. Let $\alpha_a\geq 1$ denote the redundancy
margin associated with action $a$. The formation constraints are
\begin{equation}
\label{eq:team_formation_constraints_ch4}
\beta_k(a)
\leq
\sum_{i\in\mathcal{N}} b_{ik}\mathds{1}\big(a\in\mathcal{A}_i\big)
\leq
\alpha_a\,\beta_k(a),
\quad
\sum_{k\in\mathcal{K}} b_{ik}\leq 1,
\end{equation}
where the first inequality enforces the minimum and maximum staffing level of
team $k$ for action $a$, and the second inequality guarantees disjoint team
membership.

The lower bound in~\eqref{eq:team_formation_constraints_ch4} guarantees that a
team can execute every task in its abstract plan. The upper bound prevents
excessive concentration of resources in one team. This redundancy window is
especially useful when some tasks contain partially revealed subtasks or when
travel-time estimates are conservative. The abstraction remains much smaller
than direct robot-to-subtask assignment, because the binary variables appear at
the level of teams rather than at the level of all subtasks.

Let $S_k^1$ denote the first region in $\Gamma_k$. For each robot $i$ and team
$k$, define the estimated arrival time
$t_{ik}^{\mathrm{arr}}\triangleq
\widehat{t}_i+T_{\mathrm{nav}}(\widehat{x}_i,S_k^1)$, where $\widehat{t}_i$
and $\widehat{x}_i$ are the current availability time and current position of
robot $i$. The completion cost of team $k$ is then modeled by
\begin{equation}
\label{eq:team_completion_cost_ch4}
J_k \triangleq
\max_{i\,:\,b_{ik}=1} t_{ik}^{\mathrm{arr}}
+
\sum_{\omega\in\Gamma_k} T_{\mathrm{exec}}(\omega),
\end{equation}
where the first term captures the synchronization delay of the slowest assigned
robot and the second term is the total service time of the team plan.

The team-formation layer minimizes the worst team completion time, namely
$\min \max_{k\in\mathcal{K}} J_k$, subject to
\eqref{eq:team_formation_constraints_ch4}. Once the binary solution is
obtained, the concrete team is recovered by
$\mathcal{N}_k\triangleq\{i\in\mathcal{N}\mid b_{ik}=1\}$. The outcome of this
layer is therefore a physically realizable partition of the fleet that remains
fully consistent with the abstract temporal decomposition.

\subsection{Execution-level realization of team plans}
\label{subsec:execution_realization}

After concrete teams have been formed, each abstract team plan is converted into
a rollout over regions and tasks:
\begin{equation}
\label{eq:team_rollout_ch4}
\Xi_k \triangleq
(S_k^1,\omega_k^1)\cdots(S_k^{L_k},\omega_k^{L_k}),
\end{equation}
where $S_k^\ell$ is the execution region of task $\omega_k^\ell$.
Equation~\eqref{eq:team_rollout_ch4} is the interface between the upper
coordination layers and the robot-level layer. It specifies what each team must
execute and in which temporal order, but it does not prescribe detailed robot
motions.

The execution layer resolves this remaining degree of freedom inside each team.
Given $\Xi_k$ and the team $\mathcal{N}_k$, a local coordinator generates the
robot-level plans $\tau_i$ and the associated trajectories. At this level,
motion feasibility, synchronization among robots of the same team, and the
internal assignment of subtasks are handled locally. The exact local mechanism
depends on the task family. Static known subtasks lead to routing and
trajectory-optimization problems. Partially revealed subtasks require
exploration-aware insertion policies. Dynamic subtasks require coalition updates
inside the team. These mechanisms are essential for realization, yet they do not
alter the logic of the hierarchy itself. Their role is to refine the team plan,
not to redefine it.

This separation yields the main structural advantage of the proposed framework.
Temporal satisfaction is decided at the fleet level through automaton-guided
search. Concrete resource instantiation is handled at the team-formation level
through a compact capacity-based assignment. Motion planning and local
collaboration are confined to the execution layer. Compared with monolithic
temporal-task formulations, the computational burden is therefore distributed
over smaller and better structured subproblems
\cite{kantaros2020stylus,leahy2021scalable,chen2025real,
luo2025simultaneous,gosrich2025online,luo2025hulk}. This hierarchical
decomposition is the essential reason why large fleets can be coordinated under
rich temporal tasks without explicitly constructing the full synchronous product
of all robot models and all mission automata.
